% Figure: four canonical bioreactor types: stirred tank, air-lift,
% packed-bed, hollow-fibre membrane bioreactor.
\begin{figure}[htbp]
\centering
\begin{tikzpicture}[
  vessel/.style={draw=engblack, thick, fill=engblue!5},
  liquid/.style={fill=engblue!10},
  cells/.style={fill=englime!25, draw=englime!60},
  arrow/.style={-{Stealth}, thick, draw=engred},
  gas/.style={-{Stealth}, thick, draw=engblue, densely dashed},
  lab/.style={font=\scriptsize, align=center},
]

% Stirred tank
\begin{scope}[xshift=0cm, yshift=0cm]
  \draw[vessel, liquid] (-0.9, 0) -- (-0.9, 2.4) -- (0.9, 2.4) -- (0.9, 0) arc (0:-180:0.9);
  \draw[engblack, thick] (0, 2.4) -- (0, 0.4);
  \draw[engblack, thick, fill=engblack!20] (-0.55, 0.4) rectangle (0.55, 0.6);
  \draw[engblack, thick, fill=engblack!20] (-0.55, 1.1) rectangle (0.55, 1.3);
  \draw[engblack, thick] (-0.9, 1.7) -- (-0.7, 1.7);
  \draw[engblack, thick] (0.9, 1.7) -- (0.7, 1.7);
  \draw[gas] (-0.2, -0.4) -- (-0.2, 0.0);
  \draw[gas] (0.0, -0.4) -- (0.0, 0.0);
  \draw[gas] (0.2, -0.4) -- (0.2, 0.0);
  \node[lab] at (0, -0.9) {Stirred tank\\(impellers + sparger)};
  \node[lab, anchor=west] at (-1.6, 2.1) {gas out};
  \draw[arrow] (-1.0, 2.1) -- (-0.85, 2.1);
\end{scope}

% Air-lift reactor
\begin{scope}[xshift=4.5cm, yshift=0cm]
  \draw[vessel, liquid] (-0.9, 0) -- (-0.9, 2.4) -- (0.9, 2.4) -- (0.9, 0) arc (0:-180:0.9);
  \draw[engblack, thick] (-0.25, 0.3) -- (-0.25, 2.0);
  \draw[engblack, thick] (0.25, 0.3) -- (0.25, 2.0);
  \draw[gas] (0.0, -0.4) -- (0.0, 0.0);
  \draw[arrow] (0.0, 0.5) -- (0.0, 1.9);
  \draw[arrow] (-0.55, 1.9) -- (-0.55, 0.5);
  \draw[arrow] (0.55, 1.9) -- (0.55, 0.5);
  \node[lab] at (0, -0.9) {Air-lift\\(draft tube)};
\end{scope}

% Packed bed
\begin{scope}[xshift=9.0cm, yshift=0cm]
  \draw[vessel, liquid] (-0.7, 0) rectangle (0.7, 2.6);
  \foreach \x/\y in {-0.4/0.4, 0.0/0.5, 0.4/0.4, -0.3/0.9, 0.2/1.0, -0.5/1.4, 0.0/1.5, 0.5/1.4, -0.2/1.9, 0.3/2.0, -0.45/2.3, 0.45/2.3}
    \fill[cells] (\x, \y) circle (0.18);
  \draw[arrow] (0.0, -0.4) -- (0.0, 0.0);
  \draw[arrow] (0.0, 2.6) -- (0.0, 3.0);
  \node[lab] at (0, -0.9) {Packed bed\\(immobilised cells)};
\end{scope}

% Hollow-fibre / membrane
\begin{scope}[xshift=13.5cm, yshift=0cm]
  \draw[vessel] (-1.0, 0) rectangle (1.0, 2.4);
  \foreach \y in {0.4, 0.8, 1.2, 1.6, 2.0}
    \draw[engred, thick] (-0.95, \y) -- (0.95, \y);
  \foreach \y in {0.4, 0.8, 1.2, 1.6, 2.0}
    \draw[engred, very thin] (-0.95, \y-0.08) -- (0.95, \y-0.08);
  \node[lab, anchor=west] at (1.1, 1.2) {fibre\\lumen};
  \draw[arrow] (-1.4, 1.2) -- (-1.05, 1.2);
  \draw[arrow] (1.05, 1.2) -- (1.4, 1.2);
  \draw[gas] (0, 2.4) -- (0, 2.9);
  \node[lab] at (0, -0.9) {Hollow-fibre membrane\\(perfusion / retention)};
\end{scope}

\end{tikzpicture}
\caption{Four canonical bioreactor types. The stirred-tank reactor (left) uses impellers and a sparger to disperse gas; it is the industrial default for microbial and many mammalian processes. The air-lift reactor uses a draft tube and gas injection to drive low-shear circulation. The packed-bed reactor immobilises cells or enzymes on a solid support and passes medium through the voids. The hollow-fibre membrane bioreactor retains cells in a fibre-bounded extracapillary space while medium and product cross a perfusion membrane; it supports high-density mammalian perfusion culture. The choice among the four is a transport problem (mass transfer, shear, cell retention), not a catalogue preference.}
\label{fig:vol06ch10:bioreactor-types}
\end{figure}
